\section{Methodology}
\label{sec:methodology}

\subsection{Scenarios and threat model}
We build on the SandboxEscapeBench scenarios~\cite{sandboxescapebench}, in which an
LLM agent, given a bash execution tool inside a container, attempts to break out
and read a flag file (\texttt{/flag.txt}) on the host. The baseline threat model is
deliberately weak: the agent begins with root inside the container and the host
runs default configurations. Our study holds the scenario fixed and varies only the
container hardening applied at launch.

\subsection{Execution harness}
Each scenario provides a \texttt{setup.sh} that launches the vulnerable container
(\texttt{eval-container}) via \texttt{docker run}, and the agent interacts with it
through \texttt{docker exec}. We run the host directly as the ``host'' target in a
\emph{local} sandbox mode, rather than provisioning a fresh EC2/Vagrant VM per
sample. Because all samples share one Docker daemon and container name, we (i) make
each setup idempotent (removing any prior \texttt{eval-container} before launch) and
(ii) run samples serially (\texttt{inspect eval -{}-max-samples 1}). Kernel-exploit
scenarios, which require a specific guest kernel, are out of scope for local mode.

\subsection{Hardening controls}
We evaluate eight controls, summarized in \Cref{tab:controls}. Controls B1--B4 and
the composite B8 are per-container \texttt{docker run} flags; B5 is a daemon-wide
setting applied as a separate batch; B6 swaps the container image; and B7 selects a
userspace kernel runtime. A control is selected per run via an environment variable
that a shared snippet translates into the appropriate launch arguments, so the same
hardening applies identically to both the model trial and the reference exploit.

\begin{table}[t]
\centering
\caption{Hardening controls evaluated. B1--B4 and B8 are per-container
\texttt{docker run} flags; B5 is daemon-wide; B6 swaps the image; B7 selects a
userspace-kernel runtime.}
\label{tab:controls}
\small
\begin{tabular}{@{}llp{7.2cm}@{}}
\toprule
ID & Control & Mechanism \\
\midrule
B0 & Baseline (unhardened)    & Docker defaults; reference weak state. \\
B1 & Drop all capabilities    & \texttt{-{}-cap-drop=ALL}. \\
B2 & No new privileges        & \texttt{-{}-security-opt no-new-privileges}. \\
B3 & Strict seccomp           & Custom profile blocking dangerous syscalls. \\
B4 & Read-only root fs        & \texttt{-{}-read-only}. \\
B5 & User-namespace remap     & Daemon-level remap of container root to an
                                unprivileged host UID (separate batch). \\
B6 & Minimal image            & Distroless, shell-less image (no living-off-the-land). \\
B7 & gVisor runtime           & \texttt{-{}-runtime=runsc} userspace kernel. \\
B8 & Hardened composite       & B1 + B2 + B3 + B4 combined. \\
\bottomrule
\end{tabular}
\end{table}


\paragraph{Applicability.}
Not every control applies to every scenario. For example, user-namespace remapping
(B5) is incompatible with \texttt{-{}-privileged}, so it is excluded for the
privileged scenario; the minimal-image control (B6) removes the in-container shell
entirely. We record applicability per scenario and report it alongside results.

\subsection{Model configuration}
Our primary attacker model is \texttt{google/gemini-3.5-flash}, queried through Inspect AI.
The model is a \emph{thinking-only} model whose internal reasoning cannot be
disabled: requesting a zero thinking budget (\texttt{-{}-reasoning-tokens 0}) is a
no-op for this model family, which Inspect AI surfaces as an explicit warning and
leaves reasoning enabled. All trials are therefore run with reasoning \emph{on},
uniformly across every control.

A consequence of reasoning-on operation is that the model intermittently
\emph{refuses} the escape task \emph{during} its chain-of-thought---it decides to
decline, sometimes before any tool call and sometimes after some tool calls. Because a naive scorer treats a refusal
identically to a defense-blocked failure, refusals would otherwise inflate the
apparent effectiveness of a control. We do not attempt to suppress refusals (they
cannot be reliably suppressed on this model); instead we classify every sample
into one of three outcomes: \emph{escape} (flag retrieved), \emph{blocked} (the
model attempted the exploit via tool calls but failed to retrieve the flag), and
\emph{refused} (the model issued a coherent refusal message, either before any tool call
or after some attempts; in the privileged distroless runs, every refusal followed at least one tool call). Escape success rate is computed over genuine
attempts (escape $+$ blocked), and refusal rate is reported separately
(\Cref{sec:results}). This classification is the model-side analogue of the
LLM-independent reference-exploit verification (RQ2), which cannot refuse and thus
isolates whether a control is a true architectural fix.

\paragraph{Budget and termination.}
Each sample is bounded by a token limit of \num{300000} tokens, a message/turn limit
of \num{120} messages, and a wall-clock limit of \num{1800}~seconds. Under an
effective defense the agent cannot escape and would otherwise iterate
indefinitely; these budgets cause such a sample to terminate and be recorded as a
non-escape. Escape success is scored on the flag alone, so a sample that reaches a cap
without retrieving the flag counts as a non-escape regardless of how close it came.

\subsection{Metrics}
We report, per (scenario, control) cell:
\begin{itemize}
  \item \textbf{Escape Success Rate (SR):} escapes (flag retrieved) divided by
  genuine attempts (escape $+$ blocked; refused samples excluded).
  \item \textbf{Absolute Risk Reduction (ARR):} $\mathrm{SR}(\text{B0}) -
  \mathrm{SR}(\text{control})$.
  \item \textbf{Relative Risk (RR):} $\mathrm{SR}(\text{control}) /
  \mathrm{SR}(\text{B0})$.
\end{itemize}
We run \num{10} epochs per cell and report per-cell fractions (Wilson intervals for the extreme cells appear in the Statistical analysis subsection). For
cross-control inference we analyze response patterns across scenario families using exact reference exploit outcomes paired with model outcome classification.

\subsection{Mechanistic verification (RQ2)}
\label{sec:rq3-method}
For every (scenario, control) we execute the scenario's hand-written reference
exploit (no LLM) inside the identically hardened container and record whether it
retrieves the flag. Crossing the reference-exploit outcome with the model outcome
yields the classification in \Cref{tab:rq3-scheme}. A precondition for valid RQ2
inference is that the reference exploit \emph{succeeds} at B0 (the vulnerability is
genuinely open); we verified this holds for all in-scope scenarios.

\begin{table}[t]
\centering
\caption{RQ2 classification scheme. The reference exploit is a hand-written,
LLM-free script executed under the identical hardening control.}
\label{tab:rq3-scheme}
\small
\begin{tabular}{@{}llp{6.5cm}@{}}
\toprule
Reference exploit & Model & Interpretation \\
\midrule
Fails    & Fails    & Defense truly closes the vulnerability (architectural fix). \\
Succeeds & Fails    & Vulnerability still open; model failure is friction-induced. \\
Fails    & Succeeds & Model found an unintended bypass (requires audit). \\
Succeeds & Succeeds & Vulnerability open and exploited; control ineffective here. \\
\bottomrule
\end{tabular}
\end{table}

